\begin{tabular}{lrrrrrr}
\toprule
\textbf{Det.} & \textbf{1550 V} & \textbf{1600 V} & \textbf{1650 V} & \textbf{1700 V} & \textbf{$D/D_A$} & \textbf{$V_{eq}$} \\
\midrule
\midrule
\multicolumn{7}{l}{\textbf{Cs-137}} \\
A & -- & 1.3563 & 1.8710 & 2.5502 & 1.000 & 1600 \\
B & -- & 1.0717 & 1.4364 & 1.9140 & 0.790 & 1640 \\
C & -- & 1.2055 & 1.5989 & 2.1246 & 0.889 & 1620 \\
D & -- & 0.4934 & 0.6323 & 0.7877 & 0.364 & 1824 \\
\midrule
\multicolumn{7}{l}{\textbf{Co-60}} \\
A & 0.9558 & 1.2617 & 1.6958 & 2.2694 & 1.000 & 1600 \\
B & -- & 1.0033 & 1.3472 & 1.7780 & 0.795 & 1639 \\
C & -- & 1.3868 & 1.7735 & 2.3850 & 1.099 & 1583 \\
D & -- & 0.4222 & 0.5538 & 0.7863 & 0.335 & 1783 \\
\bottomrule
\end{tabular}
\par\smallskip
\begin{minipage}{0.95\linewidth}\footnotesize Las columnas de voltaje son $D$ en mV$\cdot$ns/keV a ese alto voltaje; $D/D_A$ compara con el detector A a 1600 V. $V_{eq}$ responde la pregunta practica: a que voltaje hay que poner cada detector para que tenga la ganancia que el detector A tiene a 1600 V, usando su propia ley $D \propto V^{b}$. Extrapolar fuera de 1600--1700 V es justamente eso, extrapolar.\end{minipage}
